\chapter{Конструкторский раздел}

В данном разделе представлены функциональная модель разрабатываемого метода, требования к разрабатываемым методу и программному обеспечению, а также описание их архитектур.

\section{Функциональная модель разрабатываемого метода}

Функциональная модель разрабатываемого метода, представлена на рисунке~\ref{img:method}.

\includeimage
{method}
{f} % Обтекание (без обтекания)
{H}
{\textwidth} % Ширина рисунка
{Функциональная модель разрабатываемого метода обнаружения признаков генерации на изображении с использованием нейронной сети}

\section{Требования к разрабатываемому методу}
\label{sec:method-requirements}

Разрабатываемый метод должен отвечать следующими требованиям:
\begin{itemize}
	\item точность метода на наборе данных FaceForensics должна составлять более 95\%;
	\item количество параметров метода не должно превышать $10^8$;
	\item метод должен относится к классу <<устойчивых>>.
\end{itemize}

\section{Требования к разрабатываемому программному обеспечению}
\label{sec:software-requirements}

Разрабатываемое программное обеспечение должно обладать графическим интерфейсом, обеспечивающим пользователю возможность:
\begin{itemize}
	\item загрузить изображение;
	\item получить значение апостериорной вероятности наличия признаков генерации на загруженном изображении;
	\item получить бинарное решение о наличии признаков генерации на загруженном изображении;
	\item получить значения апостериорной вероятности наличия признаков генерации на всех лицах людей на загруженном изображении;
	\item получить бинарное решение о наличии признаков генерации на всех лицах людей на загруженном изображении;
	\item изменить пороговую вероятность функции дискриминации.
\end{itemize}

\section{Известные наборы необходимых размеченных данных}

Обучение искусственной нейронной сети требует обучающего набора размеченных данных, содержащего как подлинные, так и синтезированные изображения. Далее представлены результаты анализа известных наборов необходимых размеченных данных.

\subsection{FaceForensics}

FaceForensics -- набор данных, предложенный исследовательской группой из Технического университета Мюнхена. Данный набор предназначен для задач судебной экспертизы изображений и ориентирован на обнаружение манипуляций с лицом человека. Ключевой особенностью данного набора является его масштаб: он включает в себя более полумиллиона кадров, извлеченных из 1000 видео, полученных из открытых источников~\cite{faceforensics}.

\includeimage
{faceforensics}
{f} % Обтекание (без обтекания)
{h}
{\textwidth} % Ширина рисунка
{Пример изменения изображений с использованием метода Face2Face}

Формирование поддельных изображений в данном наборе данных выполняется с использованием метода замены лица Face2Face~\cite{face2face}. Данный метод позволяет перенести мимику и выражения лица человека из исходного видео на лицо человека из целевого видео в режиме реального времени. Пример изменения изображений с использованием данного метода представлен на рисунке~\ref{img:faceforensics}.

Преимуществом набора является его значительный объём и использование единого метода подделки Face2Face. Недостатком является ограниченность одним типом подделки, снижающая обобщающую способность моделей.

\subsection{FaceForensics++}

FaceForensics++~представляет собой расширение набора данных FaceForensics, предложенное той же исследовательской группой под руководством А. Рёсслера. Набор данных включает почти 2~миллиона изображений, полученных из 1000 видео, отобранных из открытых источников~\cite{faceforensicspp}. 

Ключевым отличием FaceForensics++ от предшественника является расширение набора методов генерации поддельных изображений. Если оригинальный набор включал только поддельные изображения, созданные с помощью метода замены лица Face2Face, то в расширенной версии добавлены четыре новых подхода к изменению лица. Данные методы представлены в таблице~\ref{tab:faceforensics-methods}.

\begin{table}[h]
	\centering
	\caption{Дополнительные методы генерации поддельных изображений в наборе данных FaceForensics++}
	\label{tab:faceforensics-methods}
	\begin{tabularx}{\textwidth}{|>{\centering\arraybackslash}p{5cm}|Y|}
		\hline
		\textbf{Метод} & \textbf{Описание} \\
		\hline
		Deepfakes~\cite{deepfakes} & Метод замены лица с использованием нейронной сети \\
		\hline
		FaceSwap~\cite{faceswap} & Метод замены лица с использованием построения его трехмерной модели \\
		\hline
		FaceShifter~\cite{faceshifter} & Метод замены лица с использованием генеративно-состязательных нейронных сетей \\
		\hline
		NeuralTextures~\cite{neural-textures} & Метод замены лица с использованием генеративно-состязательных нейронных сетей \\
		\hline
	\end{tabularx}
\end{table}

Преимуществом данного набора данных является разнообразие методов генерации подделок, повышающее репрезентативность данных для обучения моделей. Недостатком является значительный перекос по расовому и гендерному признаку~\cite{ff-minus}.

\section{Трансферное обучение нейронных сетей}

\textit{Трансферное обучение} представляет собой подход в машинном обучении, при котором знания, полученные моделью в процессе решения одной задачи, используются для повышения эффективности её обучения на другой, как правило, связанной задачи~\cite{transfer-learning}. 

В контексте задач компьютерного зрения данный процесс можно разделить на два основных этапа:
\begin{itemize}
	\item на этапе \textit{предварительного обучения} модель учится решать вспомогательную задачу, получая возможность распознавать универсальные признаки;
	\item на этапе \textit{донастройки} модель адаптируется для решения целевой задачи. 
\end{itemize}

Преимуществом трансферного обучения является ускорение сходимости и повышение итоговой точности модели при ограниченном объеме целевой выборки. Недостатком может являться риск ухудшения итоговых результатов модели, если исходная и целевая задачи являются слабо связанными~\cite{transfer-learning}.

\section{Ансамблирование нейронных сетей}

\textit{Ансамблирование нейронных сетей} представляет собой подход в машинном обучении, при котором выходные сигналы множества независимо обученных ИНС комбинируются для получения итогового решения. Принцип ансамблирования основывается на идее, что коллективное решение группы моделей в среднем оказывается точнее, чем решение отдельно взятого классификатора при условии, что ошибки этих моделей некоррелированы~\cite{ensembles}.

К преимуществам ансамблирования относятся повышение точности классификации по сравнению с использованием отдельной модели, увеличение устойчивости к шумам, а также снижение риска переобучения за счёт усреднения совокупности результатов нескольких классификаторов. Недостатками ансамблевого подхода являются пропорциональное увеличение вычислительной сложности и снижение интерпретируемости, так как итоговое решение формируется совокупностью внутренних представлений нескольких ИНС~\cite{ensembles}.

\section{Архитектура разрабатываемой нейронной сети}

Архитектура разрабатываемой нейронной сети представлена в виде ансамбля сверточной ИНС и трансформера.

\subsection{Функциональная декомпозиция архитектуры}

В результате сравнительного анализа, представленного в разделе~\ref{sec:comparative-analysis}, было установлено, что сверточные нейронные сети нацелены на выделение локальных текстурных признаков, в то время как трансформеры демонстрируют способность к анализу глобального контекста изображения. Кроме того, оба рассматриваемых класса нейронных сетей обладают устойчивостью к шуму. 

Основная идея разрабатываемой архитектуры заключается в ансамблировании сверточной ИНС и трансформера, что позволит обеспечить итоговой модели как холистическое восприятие входного изображения, так и устойчивость к искажениям.

\includeimage
{nn-idef0-0}
{f} % Обтекание (без обтекания)
{h}
{\textwidth} % Ширина рисунка
{Функциональная декомпозиция архитектуры разрабатываемой нейронной сети в виде нулевого уровня диаграммы IDEF0}

Функциональная декомпозиция архитектуры разрабатываемой нейронной сети в виде нулевого и первого уровней диаграммы IDEF0 представлена на рисунках~\ref{img:nn-idef0-0}~и~\ref{img:nn-idef0-1} соответственно.

\includeimage
{nn-idef0-1}
{f} % Обтекание (без обтекания)
{H}
{\textwidth} % Ширина рисунка
{Функциональная декомпозиция архитектуры разрабатываемой нейронной сети в виде первого уровня диаграммы IDEF0}

\subsection{Формирование ансамбля нейронных сетей}

Задача для трансферного обучения, формализованная в разделе~\ref{sec:task-formalization}, является семантически связанной с \textit{задачей распознавания образов}, заключающейся в отнесении изображения к одному из заранее определенных классов~\cite{rumelhart}. Открытый набор размеченных данных ImageNet~\cite{imagenet} предназначен для проверки пригодности моделей машинного обучения для решения задачи распознавания образов~\cite{imagenet-challenge}.

Схема архитектуры разрабатываемой нейронной сети представлена на рисунке~\ref{img:nn}.

\begin{figure}[h]
	\centering
	\begin{tikzpicture}[node distance=6mm]
		
		\tikzset{
			layer/.style={
				rectangle, draw, minimum width=6.5cm,
				minimum height=0.5cm, align=center,
				font=\small, fill=#1
			},
			layer/.default=white,
			blockbox/.style={
				draw, rectangle, rounded corners,
				inner sep=6pt
			},
			arrow/.style={
				-latex, thick
			}
		}
		
		% Вход
		\node[layer=gray!10] (input){\texttt{Изображение 224$\times$224$\times$3}};
		
		% EfficientNet
		\node[layer=gray!30, below=0.8cm of input, xshift=-4cm] (efficientnet) {\texttt{EfficientNet}};
		\node[blockbox, fit=(efficientnet)] (block1) {};
		
		\node[layer=orange!20, below=1.05cm of efficientnet] (dropout1) {\texttt{Отключение нейронов 0.3}};
		\node[layer=red!20, below=0.1cm of dropout1] (seq1) {\texttt{Полносвязный 512 + GELU}};
		\node[layer=yellow!10, below=0.1cm of seq1] (batch1) {\texttt{Пакетная нормализация}};
		\node[blockbox, fit=(dropout1) (seq1) (batch1)] (block2) {};
		
		
		% ViT
		\node[layer=gray!30, below=0.8cm of input, xshift=4cm] (vit) {\texttt{Vision Transformer}};
		\node[blockbox, fit=(vit)] (block3) {};
		
		\node[layer=orange!20, below=1.05cm of vit] (dropout2) {\texttt{Отключение нейронов 0.3}};
		\node[layer=red!20, below=0.1cm of dropout2] (seq2) {\texttt{Полносвязный 512 + GELU}};
		\node[layer=yellow!10, below=0.1cm of seq2] (batch2) {\texttt{Пакетная нормализация}};
		\node[blockbox, fit=(dropout2) (seq2) (batch2)] (block4) {};
		
		% Perceptron
		\node[layer=red!20, below=1.05cm of batch1, xshift=4cm] (seq3) {\texttt{Полносвязный 1024 + GELU}};
		\node[layer=orange!20, below=0.1cm of seq3] (dropout3) {\texttt{Отключение нейронов 0.3}};
		\node[layer=red!20, below=0.1cm of dropout3] (seq4) {\texttt{Полносвязный 256 + GELU}};
		\node[layer=red!20, below=0.1cm of seq4] (seq5) {\texttt{Полносвязный 2 + Sigmoid}};
		\node[blockbox, fit=(dropout3) (seq3) (seq4) (seq5)] (block5) {};
		
		% Result
		\node[layer=gray!10, below=0.8cm of seq5] (output) {\texttt{Результат}};
		
		% --------------------
		
		\coordinate (center1) at ($(input.south) + (0,-0.25)$);
		\coordinate (left-center1) at ($(center1) + (-4,0)$);
		\coordinate (right-center1) at ($(center1) + (4,0)$);
		
		\coordinate (center2) at ($(block5.north) + (0,0.35)$);
		\coordinate (left-center2) at ($(center2) + (-4,0)$);
		\coordinate (right-center2) at ($(center2) + (4,0)$);
		
		\draw[arrow]
		(input.south) --
		(center1) --
		(left-center1) --
		(block1.north);
		
		\draw[arrow]
		(input.south) --
		(center1) --
		(right-center1) --
		(block3.north);
		
		\draw[arrow]
		(block1.south) --
		(block2.north);
		
		\draw[arrow]
		(block3.south) --
		(block4.north);
		
		\draw[arrow]
		(block2.south) --
		(left-center2) --
		(center2) --
		(block5.north);
		
		\draw[arrow]
		(block4.south) --
		(right-center2) --
		(center2) --
		(block5.north);
		
		\draw[arrow]
		(block5.south) --
		(output.north);
		
	\end{tikzpicture}
	\caption{Схема архитектуры разрабатываемой нейронной сети}
	\label{img:nn}
\end{figure}

В качестве сверточной нейронной сети разрабатываемого ансамбля выбран EfficientNet~\cite{efficientnet} по следующим причинам:
\begin{itemize}
	\item данная нейронная сеть является самой точной известной сверточной ИНС на наборе размеченных данных ImageNet~\cite{efficientnet};
	\item значения параметров данной предобученной на наборе размеченных данных ImageNet нейронной сети находятся в открытом доступе~\cite{efficientnet-weights}.
\end{itemize}

В качестве трансформера разрабатываемого ансамбля выбран Vision Transformer~\cite{vit} по следующим причинам:
\begin{itemize}
	\item данная нейронная сеть является одним из самых точных известных трансформеров на наборе размеченных данных ImageNet~\cite{vit};
	\item значения параметров данной предобученной на наборе размеченных данных ImageNet нейронной сети находятся в открытом доступе~\cite{vit-weights}.
\end{itemize}


\subsection{Обучение разрабатываемой нейронной сети}
\label{sec:training}

В качестве набора данных для обучения разрабатываемой нейронной сети выбран FaceForensics++, содержащий 1000 подлинных и 5000 синтезированных видео с людьми. В обучающую выборку включены первые 24 кадра 1000 подлинных и 1000 синтезированных видео. Из каждого кадра предварительно извлекается лицо изображенного человека с помощью \textit{алгоритма Виолы-Джонса}~\cite{viola-jones}. Структура обучающей выборки представлена в таблицах~\ref{tab:training-sample}~и~\ref{tab:ff-sample}.

\begin{table}[h]
	\centering
	\caption{Разделение обучающей выборки разрабатываемой нейронной сети}
	\label{tab:training-sample}
	\begin{tabularx}{\textwidth}{|Y|Y|Y|Y|}
		\hline
		\textbf{Часть выборки} & \textbf{Количество видео, шт.} & \textbf{Количество изображений, шт.} & \textbf{Доля выборки, \%} \\
		\hline
		Тренировочная & 1\,800 & 43\,200 & 90 \\
		\hline
		Валидационная & 200 & 4\,800 & 10 \\
		\hline
	\end{tabularx}
\end{table}  

\begin{table}[h]
	\centering
	\caption{Разделение методов FaceForensics++ в обучающей выборке разрабатываемой нейронной сети}
	\label{tab:ff-sample}
	\begin{tabularx}{\textwidth}{|Y|Y|Y|Y|}
		\hline
		\textbf{Метод} & \textbf{Количество видео, шт.} & \textbf{Количество изображений, шт.} & \textbf{Доля выборки, \%} \\
		\hline
		Подлинный & 1000 & 24\,000 & 50 \\
		\hline
		Face2Face & 200 & 4\,800 & 10 \\
		\hline
		Deepfakes & 200 & 4\,800 & 10 \\
		\hline
		FaceSwap & 200 & 4\,800 & 10 \\
		\hline
		FaceShifter & 200 & 4\,800 & 10 \\
		\hline
		NeuralTextures & 200 & 4\,800 & 10 \\
		\hline
	\end{tabularx}
\end{table}

В процессе обучения к изображениям применяются различные виды \textit{аугментации}~\cite{augmenation}:
\begin{itemize}
	\item горизонтальное отражение;
	\item поворот на угол до 15 градусов;
	\item изменение яркости и контраста;
	\item нормализация.
\end{itemize}

Для ускорения работы алгоритма обратного распространения ошибки и защиты от переобучения используется \textit{метод AdamW}~\cite{adamw}. В качестве функции ошибки алгоритма обучения выбрана \textit{бинарная кросс-энтропия}, имеющая следующий вид:
\begin{equation}
	l(x,y) = -[y \cdot \log(\sigma(x)) + (1 - y) \cdot \log(1 - \sigma(x))]
\end{equation}
где
\begin{itemize}
	\item $x$ -- предсказанное значение;
	\item $y$ -- истинное значение;
	\item $\sigma(x)$ -- сигмоидная функция.
\end{itemize}

Схема алгоритма формирования обучающей выборки представлена на рисунке~\ref{img:samples-algorithm}.

\includeimage
{samples-algorithm}
{f} % Обтекание (без обтекания)
{H}
{0.8\textwidth} % Ширина рисунка
{Схема алгоритма формирования обучающей выборки}

\section{Архитектура разрабатываемого программного обеспечения}

Архитектура разрабатываемого программного обеспечения организована в виде слоистой монолитной модели~\cite{software-architecture}. 

\subsection{Функциональная декомпозиция архитектуры}

Функциональная декомпозиция архитектуры разрабатываемого программного обеспечения в виде нулевого и первого уровней диаграммы IDEF0 представлена на рисунках~\ref{img:software-idef0-0}~и~\ref{img:software-idef0-1} соответственно.

\includeimage
{software-idef0-0}
{f} % Обтекание (без обтекания)
{H}
{\textwidth} % Ширина рисунка
{Функциональная декомпозиция архитектуры разрабатываемого программного обеспечения в виде нулевого уровня диаграммы IDEF0}

\includeimage
{software-idef0-1}
{f} % Обтекание (без обтекания)
{H}
{0.96\textwidth} % Ширина рисунка
{Функциональная декомпозиция архитектуры разрабатываемого программного обеспечения в виде первого уровня диаграммы IDEF0}

\newpage

\subsection{Контекст проектируемой архитектуры}

Разрабатываемое программное обеспечение взаимодействует с пользователем, указывающим путь к изображению для проверки, и файловой системой, хранящей изображения и настройки используемой нейронной сети. Взаимодействие с разрабатываемым программным обеспечением происходит посредством графического интерфейса пользователя, который позволяет указать путь к изображению для проверки и отображает результат проведенного впоследствии анализа. Диаграмма контекста проектируемой архитектуры в нотации моделирования программных систем C4~\cite{c4} представлена на рисунке~\ref{img:context}.

\includeimage
{context}
{f} % Обтекание (без обтекания)
{H}
{\textwidth} % Ширина рисунка
{Диаграмма контекста проектируемой архитектуры в нотации моделирования программных систем C4}

Поток данных организован в соответствии с шаблоном проектирования Model-View-Presenter~\cite{patterns}: графический интерфейс пользователя передает команды контроллеру, который инициирует обработку команды сервисом бизнес-логики. Сервисы бизнес-логики отвечают за загрузку изображений из файловой системы, извлечение лиц из изображений и обнаружение признаков генерации на изображениях. Диаграмма потока данных проектируемой архитектуры представлена на рисунке~\ref{img:data-flow}.

\includeimage
{data-flow}
{f} % Обтекание (без обтекания)
{H}
{\textwidth} % Ширина рисунка
{Диаграмма потока данных}

\newpage

\subsection{Компоненты проектируемой архитектуры}

Диаграмма контейнеров проектируемой архитектуры в нотации моделирования программных систем C4 представлена на рисунке~\ref{img:context}.

\includeimage
{containers}
{f} % Обтекание (без обтекания)
{H}
{\textwidth} % Ширина рисунка
{Диаграмма контейнеров проектируемой архитектуры в нотации моделирования программных систем C4}

В рамках проектируемой системы выделены клиентская часть с графическим интерфейсом, серверная часть и файловое хранилище. В текущей реализации данные \textit{контейнеры}~\cite{c4} объединены в монолитное приложение. Однако выделение их в качестве самостоятельных логических единиц обеспечивает возможность последующего развертывания в виде распределенной системы без изменения глобальной логики.


\newpage

Диаграмма компонентов на языке графического описания программных систем UML представлена на рисунке~\ref{img:components}.

\includeimage
{components}
{f} % Обтекание (без обтекания)
{H}
{\textwidth} % Ширина рисунка
{Диаграмма компонентов проектируемой архитектуры на языке графического описания программных систем UML}

В рамках слоя бизнес-логики были выделены сервисы обработки и загрузки изображений, извлечения лиц и обнаружения признаков генерации на изображениях. Слой доступа к данным содержит реализации шаблона проектирования Repository~\cite{patterns} для формирования абстракции над файловой системой.

\newpage

Диаграмма классов на языке графического описания программных систем UML представлена на рисунке~\ref{img:classes}.

\includeimage
{classes}
{f} % Обтекание (без обтекания)
{H}
{\textwidth} % Ширина рисунка
{Диаграмма классов проектируемой архитектуры на языке графического описания программных систем UML}

Графический интерфейс пользователя разделен на панель управления, позволяющую загружать изображение и изменять настройки системы, и панель отображения, позволяющую производить предварительный просмотр загруженного изображения как целиком, так и отдельных содержащихся на нем лиц, а также демонстрирующей результаты проверки. Макет графического интерфейса проектируемого программного обеспечения представлен на рисунке~\ref{img:gui}.

\includeimage
{gui}
{f} % Обтекание (без обтекания)
{H}
{\textwidth} % Ширина рисунка
{Макет графического интерфейса проектируемого программного обеспечения}

\section{Выводы из конструкторского раздела}

В результате проектирования метода обнаружения признаков генерации на изображении с использованием нейронной сети были решены следующие задачи:
\begin{itemize}
	\item сформирована функциональная модель разрабатываемого метода;
	\item сформированы требования к разрабатываемому методу;
	\item сформированы требования к разрабатываемому программному обеспечению;
	\item спроектирована архитектура разрабатываемой нейронной сети;
	\item спроектированы архитектура и графический интерфейс разрабатываемого программного обеспечения.
\end{itemize}